\subjudul{Merangkum perjalanan dari class pertama hingga aplikasi utuh.}
\chapter{Penutup}

\begin{dibabini}
  \item Rangkuman materi
  \item Langkah selanjutnya
\end{dibabini}

\section{Rangkuman Materi}

Kita memulai dengan class, object, attribute, method, dan constructor untuk menyatukan data serta perilaku yang berkaitan. Encapsulation dan property menjaga perubahan data agar mengikuti aturan. Class method, static method, dan decorator memperluas cara menyusun perilaku pada class.

Inheritance, abstraksi, dan polymorphism membantu membentuk kontrak serta variasi perilaku. Association, aggregation, dan composition menjelaskan hubungan antar-object. Pemilihannya mengikuti kebutuhan serta tanggung jawab class.

Module dan package mengatur kode, exception handling menangani kegagalan, dan file handling mempertahankan data setelah program ditutup. Melalui GUI, input pengguna dihubungkan dengan operasi aplikasi menggunakan widget dan event handler.

Proyek perpustakaan mempertemukan materi tersebut: menganalisis kebutuhan, merancang class dan relasi, mengimplementasikan aturan, lalu membuat antarmuka. Kita belajar menjaga data, menjelaskan penolakan, serta memeriksa hasil setelah aplikasi disimpan dan dibuka kembali.

\section{Langkah Selanjutnya}

Kemampuan merancang program akan berkembang melalui latihan dan perbaikan. Pilih satu pengembangan kecil dari proyek, misalnya pencarian buku atau laporan peminjaman terlambat. Tuliskan hasil yang diharapkan, tentukan class yang bertanggung jawab, lalu implementasikan dan periksa skenario berhasil maupun gagal.

Cobalah pula menerapkan langkah yang sama pada masalah lain, seperti peminjaman alat laboratorium atau pengelolaan stok sederhana. Gunakan dokumentasi Python ketika menemui perilaku yang belum dipahami. Saat terjadi kesalahan, baca pesannya, telusuri perubahan data, dan perbaiki penyebabnya satu per satu.

\begin{tips}[Terus Mencoba dan Berlatih]
Terima kasih telah belajar bersama melalui buku ini. Luangkan waktu untuk menulis kode sendiri, mengubah contoh, dan memeriksa hasilnya. Setiap latihan memberi kesempatan untuk memahami keputusan rancangan dengan lebih baik. Selamat mencoba dan terus berlatih agar semakin mahir menggunakan OOP.
\end{tips}
